\SourceLeaf{329}

\begin{center}
[4.]

[\emph{Handbuch}, S.\,203.]
\end{center}

Die Auflösung der \textsc{Pothenot}schen Aufgabe beruht auf den drei Gleichungen
\[
\frac{\alpha a+\beta b+\gamma c}{\alpha}\, . \, x' - \frac{\alpha' a'+\beta' b'+\gamma' c'}{\alpha'}\, . \, x - \frac{\beta b+\gamma c}{\alpha}\, . \, a' + \frac{\beta' b'+\gamma' c'}{\alpha'}\, . \, a = 0
\]
\[
\frac{\alpha a+\beta b+\gamma c}{\beta}\, . \, x' - \frac{\alpha' a'+\beta' b'+\gamma' c'}{\beta'}\, . \, x - \frac{\alpha a+\gamma c}{\beta}\, . \, b' + \frac{\alpha' a'+\gamma' c'}{\beta'}\, . \, b = 0
\]
\[
\frac{\alpha a+\beta b+\gamma c}{\gamma}\, . \, x' - \frac{\alpha' a'+\beta' b'+\gamma' c'}{\gamma'}\, . \, x - \frac{\alpha a+\beta b}{\gamma}\, . \, c' + \frac{\alpha' a'+\beta' b'}{\gamma'}\, . \, c = 0,
\]

wovon jede schon in den beiden andern enthalten ist. Also
\[
a\alpha+b\beta+c\gamma=\omega \quad [\text{und } a'\alpha'+b'\beta'+c'\gamma'=\omega']
\]
gesetzt:
\[
\left(\frac{\beta}{\beta'}-\frac{\gamma}{\gamma'}\right)\omega' x + (b'-c')\omega - \left(\frac{\beta b}{\beta'}-\frac{\gamma c}{\gamma'}\right)\omega' = 0
\]
\[
\left(\frac{\gamma}{\gamma'}-\frac{\alpha}{\alpha'}\right)\omega' x + (c'-a')\omega - \left(\frac{\gamma c}{\gamma'}-\frac{\alpha a}{\alpha'}\right)\omega' = 0
\]
\[
\left(\frac{\alpha}{\alpha'}-\frac{\beta}{\beta'}\right)\omega' x + (a'-b')\omega - \left(\frac{\alpha a}{\alpha'}-\frac{\beta b}{\beta'}\right)\omega' = 0.
\]

Noch zierlicher
\[
\begin{aligned}
&(\alpha\beta'-\beta\alpha')(\alpha' a'+\beta' b'+\gamma' c')x + a\alpha\beta'\gamma'(b'-c')\\
&\quad+ \beta b\alpha'\gamma'(c'-a') + \gamma c\beta'\alpha'(a'-b') = 0.
\end{aligned}
\]

\medskip

Sind $A, B, C$ drei complexe Zahlen, deren Winkel den beobachteten Azimuthen entsprechen, so kann man setzen
\[
\alpha = A(BC'-CB')
\]
\[
\beta = B(CA'-AC')
\]
\[
\gamma = C(AB'-BA').
\]

Die physische Möglichkeit beruht darauf, dass
\[
B'C(A-B)(A'-C') - BC'(A'-B')(A-C)
\]
mit den beiden andern symmetrisch daraus folgenden Grössen einerlei Zeichen haben [muss].
